%% Andrea Piccolo - CV
%% Tailored for: Technical Software Engineer, S[&]T High-Tech Consulting, Delft
%%
%% Compile with: lualatex andrea_piccolo_cv_stcorp_technical-software-engineer.tex
%% (pdflatex often fails on modern MiKTeX with fontawesome5 font-expansion errors.)

\documentclass[11pt,a4paper,sans]{moderncv}
\moderncvstyle{banking}
\moderncvcolor{blue}

% Force both first and last name AND section headings to render in moderncv
% blue (color1). Default banking on lualatex+MiKTeX leaves these black, which
% looks inconsistent with the rest of the blue accent scheme.
\renewcommand*{\firstnamestyle}[1]{{\fontsize{34}{36}\bfseries\upshape\color{color1}#1}}
\renewcommand*{\lastnamestyle}[1]{{\fontsize{34}{36}\bfseries\upshape\color{color1}#1}}
\renewcommand*{\sectionstyle}[1]{{\sectionfont\color{color1}#1}}

\usepackage[utf8]{inputenc}
\usepackage{hyperref}
\hypersetup{
    colorlinks=true,
    linkcolor=blue,
    filecolor=magenta,
    urlcolor=blue,
    pdftitle={Andrea Piccolo - CV},
    pdfpagemode=FullScreen,
}
\usepackage[scale=0.77]{geometry}
\usepackage{needspace}
\usepackage{import}
\hyphenation{FEniCSx StonedFEniCSx}

% personal data
\name{Andrea}{Piccolo}
\address{Leeds, UK}{}{}
\phone[mobile]{+44 7831 007579}
\email{andreapiccolope@gmail.com}
\extrainfo{\href{https://www.linkedin.com/in/ap0889/}{LinkedIn}, \href{https://github.com/APiccolo89}{GitHub} \textbar{} Italian citizen -- EU work rights, no permit needed}

\begin{document}

\makecvtitle

% ============================================================
%     PROFILE STATEMENT
% ============================================================

\vspace{6pt}
\small{Scientific software engineer with a PhD in Computational Geodynamics and ten years writing numerical software in Python, Julia and C on Linux and HPC clusters. Builds tested, documented packages from research prototypes, most recently StonedFEniCSx, an open-source FEM package with CI, regression tests and Sphinx documentation. Optimised an existing scientific package by 10\texttimes{} and condensed 20 TB of simulation output into 200 GB of public databases used by other groups.}

% ============================================================
%     SKILLS
% ============================================================

\section{Skills}
\vspace{1pt}
\begin{itemize}

\item \textbf{Software Development}: Object-oriented Python (NumPy, SciPy, pandas, h5py, Numba, FEniCSx), Julia, C, MATLAB; turning research scripts into maintainable packages (class hierarchies, dataclass configuration, YAML input, pyproject packaging)

\item \textbf{Testing, Review \& Documentation}: pytest with Codecov coverage, regression tests against published benchmarks, GitHub Actions CI, pull-request workflow in shared open-source codebases, Sphinx documentation deployed to GitHub Pages

\item \textbf{Data Processing \& Optimisation}: 10\texttimes{} speedup of an existing scientific package; 100+ TB of simulation data processed; automated pipelines for running, post-processing and analysing simulations; 2D/3D visualisation tools

\item \textbf{Linux \& HPC}: Daily Linux work on HPC clusters (SLURM, PBS), MPI-parallel FEM (FEniCSx/PETSc), reproducible environments (venv/conda), Git; SQL (learning)

\item \textbf{Collaboration}: Cross-institutional projects (Mainz, Bayreuth, KIT Karlsruhe, Leeds); supplying data, tools and support to partner groups; conference presentations; supervision and teaching

\end{itemize}

% ============================================================
%     PROFESSIONAL EXPERIENCE
% ============================================================

\section{Professional Experience}
\vspace{3pt}
\begin{itemize}

\needspace{5\baselineskip}
\item{\cventry{02/2025--Present}{Research Fellow}{University of Leeds}{Leeds, UK}{}{\vspace{1pt}
\begin{itemize}
    \item Built \href{https://github.com/APiccolo89/StonedFenicsx}{StonedFEniCSx}, an open-source FEniCSx package for coupled thermal and Stokes modelling, growing it from a single script into an object-oriented, pip-installable package: problem and solver class hierarchies, dataclass configuration, YAML input
    \item Set up CI with GitHub Actions: pytest suite with Codecov coverage, regression tests against the van Keken et al.\ benchmarks, and Sphinx documentation built and deployed to GitHub Pages
    \item Optimised a pre-existing numerical package, achieving a 10\texttimes{} speedup
    \item Built automated workflows to run simulations, post-process outputs and analyse large datasets; MPI-parallel runs through FEniCSx/PETSc
\end{itemize}}}

\vspace{3pt}

\item{\cventry{11/2021--05/2024}{Research Fellow}{University of Bayreuth}{Bayreuth, Germany}{}{\vspace{1pt}
\begin{itemize}
    \item Extended an open-source Julia package, coupling forward mechanical modelling with basin modelling
    \item Delivered public databases: 20 TB of simulated data condensed into 200 GB of accessible outputs
    \item Designed a 0D reduced-order model predicting 3D simulation behaviour, plus an optimisation algorithm recovering parameter combinations that match experimental data
    \item Supplied simulation data, Python tools and support to the KIT Karlsruhe basin modelling group
\end{itemize}}}

\vspace{3pt}

\needspace{6\baselineskip}
\item{\cventry{11/2015--07/2021}{Research Fellow}{Johannes Gutenberg Universität Mainz}{Mainz, Germany}{}{\vspace{1pt}
\begin{itemize}
    \item Incorporated chemical modelling into a large-scale 3D parallel numerical code (\href{https://github.com/UniMainzGeo/LaMEM}{LaMEM}, C)
    \item Generated and processed 100+ TB of simulation data on HPC clusters; developed 2D/3D visualisation tools
    \item Presented results at international conferences; supported supervision and teaching
\end{itemize}}}

\end{itemize}

% ============================================================
%     EDUCATION
% ============================================================

\section{Education}
\vspace{1pt}
\begin{itemize}

\needspace{5\baselineskip}
\item{\cventry{11/2015--05/2019}{PhD in Computational Geodynamics}{Johannes Gutenberg Universität Mainz}{Mainz, Germany}{}{\vspace{1pt}
Thermo-mechanical numerical modelling, large-scale data analysis, CI.
}}

\vspace{3pt}

\item{\cventry{10/2012--03/2015}{MSc in Geology and Technical Geology}{University of Padua}{Padua, Italy}{}{\vspace{1pt}
Data analysis, numerical modelling, scientific communication.
}}

\vspace{3pt}

\item{\cventry{10/2008--03/2012}{BSc in Earth Sciences}{University of Padua}{Padua, Italy}{}{\vspace{1pt}
Data analysis, MATLAB.
}}

\end{itemize}

% ============================================================
%     REPOSITORIES & SELECTED PUBLICATIONS
% ============================================================

\section{Repositories}
\vspace{1pt}
\begin{itemize}
\item \textbf{Author:} \href{https://github.com/APiccolo89/StonedFenicsx}{StonedFEniCSx} \textbar{} \textbf{Contributor:} \href{https://github.com/UniMainzGeo/LaMEM}{LaMEM}, \href{https://github.com/JuliaGeodynamics/GeophysicalModelGenerator.jl}{GeophysicalModelGenerator.jl}, \href{https://bitbucket.org/bkaus/mvep2/}{MVEP2}
\end{itemize}

\needspace{5\baselineskip}
\section{Selected Publications}
\vspace{1pt}
{\small
\begin{itemize}
\item Kaus, B.\,J., \ldots, \textbf{Piccolo, A.}, et al.\ (2024). GeophysicalModelGenerator.jl: A Julia package to visualise geoscientific data and create numerical model setups. \textit{Journal of Open Source Software}, 9(103), 6763.
\item \textbf{Piccolo, A.}, Spang, A., et al.\ (2026). The dynamics and surface signal of slab break-off in continental settings: Insights from 3D numerical modeling. \textit{JGR: Solid Earth}, 131(5).
\item \textbf{Piccolo, A.}, Palin, R.\,M., et al.\ (2019). Generation of Earth's early continents from a relatively cool Archean mantle. \textit{Geochem., Geophys., Geosyst.}, 20(4), 1679--1697.
\item[] 9 further peer-reviewed papers (2018--2026); full list on request.
\end{itemize}}

% ============================================================
%     LANGUAGES
% ============================================================

\section{Languages}
\vspace{1pt}
\begin{itemize}
\item Italian (native), English (C1), German (B1).
\end{itemize}

% ============================================================
%     REFERENCES
% ============================================================

\section{References}
\vspace{1pt}
\begin{itemize}
\item Prof.\ Dr.\ Timothy J.\ Craig --- University of Leeds (\href{mailto:T.J.Craig@leeds.ac.uk}{T.J.Craig@leeds.ac.uk})
\item Prof.\ Dr.\ Marcell Thielmann --- University of Bonn (\href{mailto:thielmann@geo.uni-bonn.de}{thielmann@geo.uni-bonn.de})
\item Prof.\ Dr.\ Nevena Andri\'c Toma\v{s}evi\'c --- KIT Karlsruhe (\href{mailto:nevena.tomasevic@kit.edu}{nevena.tomasevic@kit.edu})
\item Dr.\ Arne Spang --- University of Bayreuth (\href{mailto:arne.spang@uni-bayreuth.de}{arne.spang@uni-bayreuth.de})
\end{itemize}

\end{document}
